\ifdefined\gkver\else\def\gkver{student}\fi
\documentclass[\gkver]{gaokaozhenti}
\newcommand{\ccnum}[1]{{\fontspec{Noto Sans CJK JP}#1}}
\begin{document}

\chapter{2017年全国理综Ⅰ}
%% paperType: 理综物理部分

\begin{enumerate}
\item
%% number: 14
%% paperName: 2017年全国理综Ⅰ第 14 题 6 分
%% typeId: 1
%% type: 单选题
%% chapter: 运动和力
%% point: 动量守恒定律
%% method: 无
%% score: 6
%% degree: 950
%% duplicateId: 0
%% body:
将质量为 $1.00\Ukg$ 的模型火箭点火升空，$50\Ug$ 燃烧的燃气以大小为 $600\Ums$ 的速度从火箭喷口在很短时间内喷出。在燃气喷出后的瞬间，火箭的动量大小为（喷出过程中重力和空气阻力可忽略）\xzanswer{A}

\fourchoices[answer=A]
{$30\Ukgms$}
{$5.7\times10^{2}\Ukgms$}
{$6.0\times10^{2}\Ukgms$}
{$6.3\times10^{2}\Ukgms$}

%% answer: A
%% memo:
\memoanswer{
	设火箭的总质量为 $M$，燃气的质量为 $m$，取火箭的运动方向为正方向，由动量守恒定律得 $0=(M-m)v-mv_{0}$，其中 $v_{0}=600\Ums$，则火箭的动量大小为 $p=(M-m)v=mv_{0}=50\times10^{-3}\times600\Ukgms=30\Ukgms$，故 A 正确，BCD 错误。
}

\item
%% number: 15
%% paperName: 2017年全国理综Ⅰ第 15 题 6 分
%% typeId: 1
%% type: 单选题
%% chapter: 曲线运动
%% point: 平抛运动
%% method: 无
%% score: 6
%% degree: 750
%% duplicateId: 0
%% body:
发球机从同一高度向正前方依次水平射出两个速度不同的乒乓球（忽略空气的影响）。速度较大的球越过球网，速度较小的球没有越过球网，其原因是\xzanswer{C}

\fourchoices[answer=C]
{速度较小的球下降相同距离所用的时间较多}
{速度较小的球在下降相同距离时在竖直方向上的速度较大}
{速度较大的球通过同一水平距离所用的时间较少}
{速度较大的球在相同时间间隔内下降的距离较大}

%% answer: C
%% memo:
\memoanswer{
	将乒乓球的平抛运动分解为水平方向的匀速直线运动和竖直方向的自由落体运动，由在竖直方向 $h=\frac{1}{2}gt^{2}$、$v_{y}=gt$，知在竖直方向下降相同的高度时，两球所用时间相同、竖直方向上的速度相同，故 ABD 错误；乒乓球在水平方向做匀速直线运动，有 $x=v_{0}t$，则速度较大的球通过同一水平距离所用的时间较少，故 C 正确。
}

\item
%% number: 16
%% paperName: 2017年全国理综Ⅰ第 16 题 6 分
%% typeId: 1
%% type: 单选题
%% chapter: 磁场
%% point: 带电粒子在复合场中的运动
%% method: 无
%% score: 6
%% degree: 650
%% duplicateId: 0
%% body:
如图，空间某区域存在匀强电场和匀强磁场，电场方向竖直向上（与纸面平行），磁场方向垂直于纸面向里，三个带正电的微粒 a，b，c 电荷量相等，质量分别为 $m_{a}$，$m_{b}$，$m_{c}$，已知在该区域内，a 在纸面内做匀速圆周运动，b 在纸面内向右做匀速直线运动，c 在纸面内向左做匀速直线运动。下列选项正确的是\xzanswer{B}

\onepicture[h=3.6cm]{figs/16.png}

\fourchoices[answer=B]
{$m_{a}>m_{b}>m_{c}$}
{$m_{b}>m_{a}>m_{c}$}
{$m_{c}>m_{b}>m_{a}$}
{$m_{a}>m_{c}>m_{b}$}

%% answer: B
%% memo:
\memoanswer{
	设三个微粒带电荷量为 $q$，分析微粒在叠加场中的受力可知，微粒均受到三个作用力，即重力、电场力、洛伦兹力。a 做匀速圆周运动，则说明重力与电场力平衡，a 所受的合力为洛伦兹力，提供其做圆周运动的向心力，即 $m_{a}g=qE$；b 向右做匀速直线运动，即合外力为零，结合左手定则可知 $m_{b}g=qE+qv_{b}B$；同理，c 向左做匀速直线运动，有 $m_{c}g=qE-qv_{c}B$，故 $m_{b}>m_{a}>m_{c}$，故 B 正确。
}

\item
%% number: 17
%% paperName: 2017年全国理综Ⅰ第 17 题 6 分
%% typeId: 1
%% type: 单选题
%% chapter: 原子和原子核
%% point: 质能方程
%% method: 无
%% score: 6
%% degree: 750
%% duplicateId: 0
%% body:
大科学工程“人造太阳”主要是将氘核聚变反应释放的能量用来发电，氘核聚变反应方程是 $\ce{^2_1 H + ^2_1 H -> ^3_2 He + ^1_0 n}$，已知 $\ce{^2_1 H}$ 的质量为 $2.0136\Uu$，$\ce{^3_2 He}$ 的质量为 $3.0150\Uu$，$\ce{^1_0 n}$ 的质量为 $1.0087\Uu$，$1\Uu=931\UMeV/c^{2}$。氘核聚变反应中释放的核能约为\xzanswer{B}

\fourchoices[answer=B]
{$3.7\UMeV$}
{$3.3\UMeV$}
{$2.7\UMeV$}
{$0.93\UMeV$}

%% answer: B
%% memo:
\memoanswer{
	反应过程中亏损的质量为 $\Delta m=2\times2.0136\Uu-1.0087\Uu-3.0150\Uu=0.0035\Uu$，反应过程中释放的核能 $\Delta E=0.0035\times931\UMeV\approx3.3\UMeV$，故 B 正确。
}

\item
%% number: 18
%% paperName: 2017年全国理综Ⅰ第 18 题 6 分
%% typeId: 1
%% type: 单选题
%% chapter: 电磁感应
%% point: 楞次定律推广
%% method: 无
%% score: 6
%% degree: 650
%% duplicateId: 0
%% body:
扫描隧道显微镜（STM）可用来探测样品表面原子尺寸上的形貌，为了有效隔离外界振动对 STM 的扰动，在圆底盘周边沿其径向对称地安装若干对紫铜薄板，并施加磁场来快速衰减其微小振动，如图所示，无扰动时，按下列四种方案对紫铜薄板施加恒磁场；出现扰动后，对于紫铜薄板上下及其左右振动的衰减最有效的方案是\xzanswer{A}

\onepicture[h=3.8cm]{\includesvg{figs/18}}

\fourchoices[answer=A, ispicture=true, h=1.8cm]
{figs/18a.png}
{figs/18b.png}
{figs/18c.png}
{figs/18d.png}

%% answer: A
%% memo:
\memoanswer{
	施加磁场来衰减振动是利用电磁阻尼实现的，即振动过程中发生电磁感应现象，产生的感应电流在磁场中受安培力阻碍其振动。A 图中紫铜薄板无论上下振动还是左右振动都会出现磁通量的变化，即产生电磁感应现象；B 图中紫铜薄板上下振动时没有磁通量的变化，不会产生电磁感应现象；C 图中紫铜薄板无论上下振动还是左右振动均不会产生电磁感应现象；D 图中紫铜薄板上下振动时不会产生电磁感应现象，故最有效的方案是 A，故 A 正确。
}

\item
%% number: 19
%% paperName: 2017年全国理综Ⅰ第 19 题 6 分
%% typeId: 2
%% type: 多选题
%% chapter: 磁场
%% point: 安培力
%% method: 无
%% score: 6
%% degree: 550
%% duplicateId: 0
%% body:
如图，三根相互平行的固定长直导线 $L_{1}$、$L_{2}$ 和 $L_{3}$ 两两等距，均通有电流 $I$，$L_{1}$ 中电流方向与 $L_{2}$ 中的相同，与 $L_{3}$ 中的相反，下列说法正确的是\xzanswer{BC}

\onepicture[h=4.2cm]{figs/19.png}

\fourchoices[answer=BC]
{$L_{1}$ 所受磁场作用力的方向与 $L_{2}$、$L_{3}$ 所在平面垂直}
{$L_{3}$ 所受磁场作用力的方向与 $L_{1}$、$L_{2}$ 所在平面垂直}
{$L_{1}$、$L_{2}$ 和 $L_{3}$ 单位长度所受的磁场作用力大小之比为 $1:1:\sqrt{3}$}
{$L_{1}$、$L_{2}$ 和 $L_{3}$ 单位长度所受的磁场作用力大小之比为 $\sqrt{3}:\sqrt{3}:1$}

%% answer: BC
%% memo:
\memoanswer{
	因三根导线对称放置，构成等边三角形，每根导线通有相等的电流，由安培定则可知每根导线在三角形顶点处产生的磁感应强度大小相等，均为 $B_{0}$。由磁感应强度的叠加原理可知 $L_{1}$ 处的磁感应强度方向垂直于 $L_{2}$、$L_{3}$ 所在平面斜向上，大小为 $B_{0}$，由左手定则可知 $L_{1}$ 所受磁场作用力的方向与 $L_{2}$、$L_{3}$ 所在平面平行，故 A 错误；同理，$L_{3}$ 处的合磁场的磁感应强度方向平行于 $L_{1}$、$L_{2}$ 所在平面向右，大小为 $\sqrt{3}B_{0}$，由左手定则可知 $L_{3}$ 所受磁场作用力的方向与 $L_{1}$、$L_{2}$ 所在平面垂直，故 B 正确；由对称性可知，$L_{1}$、$L_{2}$ 单位长度所受的磁场作用力大小相等，均为 $F_{1}=F_{2}=\frac{B_{0}IL}{L}=B_{0}I$，$L_{3}$ 单位长度所受的磁场作用力大小为 $F_{3}=\frac{\sqrt{3}B_{0}IL}{L}=\sqrt{3}B_{0}I$，所以三者的大小之比为 $1:1:\sqrt{3}$，故 C 正确，D 错误。
}

\item
%% number: 20
%% paperName: 2017年全国理综Ⅰ第 20 题 6 分
%% typeId: 2
%% type: 多选题
%% chapter: 电场
%% point: 电场中的图象
%% method: 无
%% score: 6
%% degree: 450
%% duplicateId: 0
%% body:
在一静止点电荷的电场中，任一点的电势 $\varphi$ 与该点到点电荷的距离 $r$ 的关系如图所示。电场中四个点 a、b、c 和 d 的电场强度大小分别 $E_{a}$、$E_{b}$、$E_{c}$ 和 $E_{d}$。点 a 到点电荷的距离 $r_{a}$ 与点 a 的电势 $\varphi_{a}$ 已在图中用坐标 $(r_{a},\varphi_{a})$ 标出，其余类推。现将一带正电的试探电荷由 a 点依次经 b、c 点移动到 d 点，在相邻两点间移动的过程中，电场力所做的功分别为 $W_{ab}$、$W_{bc}$ 和 $W_{cd}$。下列选项正确的是\xzanswer{AC}

\onepicture[h=4.4cm]{figs/20.png}

\fourchoices[answer=AC]
{$E_{a}:E_{b}=4:1$}
{$E_{c}:E_{d}=2:1$}
{$W_{ab}:W_{bc}=3:1$}
{$W_{bc}:W_{cd}=1:3$}

%% answer: AC
%% memo:
\memoanswer{
	由点电荷的电场强度表达式 $E=\frac{kQ}{r^{2}}$，则 $\frac{E_{a}}{E_{b}}=\frac{r_{b}^{2}}{r_{a}^{2}}=\frac{4}{1}$，故 A 正确；同理，$\frac{E_{c}}{E_{d}}=\frac{4}{1}$，故 B 错误；由电势能与电场力做功的关系 $W_{ab}=q\varphi_{a}-q\varphi_{b}$，$\frac{W_{ab}}{W_{bc}}=\frac{\varphi_{a}-\varphi_{b}}{\varphi_{b}-\varphi_{c}}=\frac{6-3}{3-2}=3$，故 C 正确；同理，$\frac{W_{bc}}{W_{cd}}=1$，故 D 错误。
}

\item
%% number: 21
%% paperName: 2017年全国理综Ⅰ第 21 题 6 分
%% typeId: 2
%% type: 多选题
%% chapter: 力物体的平衡
%% point: 动态平衡
%% method: 作图法
%% score: 6
%% degree: 450
%% duplicateId: 0
%% body:
如图，柔软轻绳 ON 的一端 O 固定，其中间某点 M 拴一重物，用手拉住绳的另一端 N，初始时，$OM$ 竖直且 $MN$ 被拉直，$OM$ 与 $MN$ 之间的夹角为 $\alpha$（$\alpha>\frac{\pi}{2}$）。现将重物向右上方缓慢拉起，并保持夹角 $\alpha$ 不变。在 $OM$ 由竖直被拉到水平的过程中\xzanswer{AD}

\onepicture[h=4.8cm]{\includesvg{figs/21}}

\fourchoices[answer=AD]
{MN 上的张力逐渐增大}
{MN 上的张力先增大后减小}
{OM 上的张力逐渐增大}
{OM 上的张力先增大后减小}

%% answer: AD
%% memo:
\memoanswer{
	对重物受力分析，画出甲、乙、丙三个特殊位置的受力图，其中 $T_{OM}$ 和 $T_{MN}$ 的合力大小、方向均不变，大小等于 $G$，在重物移动的过程中 $OM$ 与 $MN$ 的夹角 $\alpha$ 不变，由图甲、乙、丙可知 $T_{OM}$ 先增大后减小，$T_{MN}$ 逐渐增大，故 AD 正确。

	\onepicture[w=4.4cm]{figs/21_an.png}
}

\item
%% number: 22
%% paperName: 2017年全国理综Ⅰ第 22 题 5 分
%% typeId: 4
%% type: 实验题
%% chapter: 直线运动
%% point: 打点计时器公式
%% method: 无
%% score: 5
%% degree: 550
%% duplicateId: 0
%% body:
某探究小组为了研究小车在桌面上的直线运动，用自制“滴水计时器”计量时间。实验前，将该计时器固定在小车旁，如图\ref{2017全国理综Ⅰ22a}所示。实验时保持桌面水平，用手轻推一下小车。在小车运动过程中，滴水计时器等时间间隔地滴下小水滴，图\ref{2017全国理综Ⅰ22b}记录了桌面上连续的 6 个水滴的位置。（已知滴水计时器每 $30\Us$ 内共滴下 46 个小水滴）

\begin{enumerate}
	\item
	由图\ref{2017全国理综Ⅰ22b}可知，小车在桌面上是\tkanswer[2.2cm]{从右向左}（填“从右向左”或“从左向右”）运动的。

	\item
	该小组同学根据图\ref{2017全国理综Ⅰ22b}的数据判断出小车做匀变速运动。小车运动到图\ref{2017全国理综Ⅰ22b}中 A 点位置时的速度大小为\tkanswer[2.0cm]{0.19}$\Ums$，加速度大小为\tkanswer[2.0cm]{0.037}$\Umsq$。（结果均保留 2 位有效数字）
\end{enumerate}

\twopicture[numA={（a）}, numB={（b）}, labelA=2017全国理综Ⅰ22a, labelB=2017全国理综Ⅰ22b, hA=4.0cm, hB=3.0cm]
{figs/22a.png}
{figs/22b.png}

%% answer:
\jdanswer{
\begin{enumerate}
	\item
	从右向左

	\item
	$0.19\Ums$，$0.037\Umsq$
\end{enumerate}
}
%% memo:
\memoanswer{
	（1）小车受到摩擦力作用做减速运动，水滴间距离逐渐减小，故小车从右向左运动。

	（2）滴水的时间间隔 $T=\frac{30}{45}\Us=\frac{2}{3}\Us$，小车运动到 A 点时的速度 $v=\frac{0.117+0.133}{2\times\frac{2}{3}}\Ums=0.19\Ums$，小车的加速度 $a=\frac{(0.150+0.133)-(0.117+0.100)}{4\times(\frac{2}{3})^{2}}\Umsq=0.037\Umsq$。
}

\item
%% number: 23
%% paperName: 2017年全国理综Ⅰ第 23 题 10 分
%% typeId: 4
%% type: 实验题
%% chapter: 电路
%% point: 伏安法测电阻
%% method: 无
%% score: 10
%% degree: 450
%% duplicateId: 0
%% body:
某同学研究小灯泡的伏安特性，所使用的器材有：小灯泡 L（额定电压 $3.8\UV$，额定电流 $0.32\UA$）；电压表 V（量程 $3\UV$，内阻 $3\UkO$）；电流表 A（量程 $0.5\UA$，内阻 $0.5\UO$）；固定电阻 $R_{0}$（阻值 $1000\UO$）；滑动变阻器 $R$（阻值 $0\sim9.0\UO$）；电源 $E$（电动势 $5\UV$，内阻不计）；开关 S；导线若干。

\begin{enumerate}
	\item
	实验要求能够实现在 $0\sim3.8\UV$ 的范围内对小灯泡的电压进行测量，画出实验电路原理图。

	\item
	实验测得该小灯泡伏安特性曲线如图\ref{2017全国理综Ⅰ23a}所示。由实验曲线可知，随着电流的增加小灯泡的电阻\tkanswer[2.0cm]{增大}（填“增大”、“不变”或“减小”），灯丝的电阻率\tkanswer[2.0cm]{增大}（填“增大”、“不变”或“减小”）。

	\item
	用另一电源 $E_{0}$（电动势 $4\UV$，内阻 $1.00\UO$）和题给器材连接成图\ref{2017全国理综Ⅰ23b}所示的电路，调节滑动变阻器 $R$ 的阻值，可以改变小灯泡的实际功率。闭合开关 S，在 $R$ 的变化范围内，小灯泡的最小功率为\tkanswer[2.0cm]{0.39}$\UW$，最大功率为\tkanswer[2.0cm]{1.17}$\UW$。（结果均保留 2 位小数）
\end{enumerate}

\twopicture[numA={（a）}, numB={（b）}, labelA=2017全国理综Ⅰ23a, labelB=2017全国理综Ⅰ23b, hA=4.0cm, hB=3.4cm]
{figs/23a.png}
{figs/23b.png}

%% answer:
\jdanswer{
\begin{enumerate}
	\item
	如图所示

	\onepicture[h=3.4cm]{figs/23_an1.png}

	\item
	增大；增大

	\item
	$0.39\UW$，$1.17\UW$
\end{enumerate}
}
%% memo:
\memoanswer{
	（1）实验要求能够实现在 $0\sim3.8\UV$ 的范围内进行测量，电压表量程为 $3\UV$，需要改装电压表增大量程，故电压表应串联定值电阻，滑动变阻器需要连接成分压式电路，由于 $R_{x}<\sqrt{R_{A}R_{V}}$，故电流表应采用外接法。

	（2）伏安特性曲线上的点和原点连线的斜率表示电阻的倒数，图中曲线上的点和原点连线的斜率逐渐减小，则电阻逐渐增大；电流增大时，灯丝的温度升高，金属材料的电阻率随着温度的升高而增大。

	（3）利用作图法，将滑动变阻器等效为电源内阻，当电源等效内阻为 $r'=r+R=10\UO$ 时，小灯泡消耗的功率最小，这时作出电源 $I-U$ 的图像如图\ccnum{①}所示（此时斜率为 0.1），可知小灯泡消耗最小功率为 $P_{\min}=U_{1}I_{1}=1.75\times0.22\UW=0.39\UW$；当电源等效内阻为 $r'=r+0=1\UO$ 时，小灯泡消耗的功率最大，这时作出电源 $I-U$ 的图像如图\ccnum{②}所示（此时斜率为 1），可知小灯泡消耗最大功率为 $P_{\max}=3.67\times0.318\UW=1.17\UW$。

	\onepicture[w=6.6cm]{figs/23_an2.png}
}

\item
%% number: 24
%% paperName: 2017年全国理综Ⅰ第 24 题 12 分
%% typeId: 6
%% type: 计算题
%% chapter: 功和能
%% point: 功能原理
%% method: 无
%% score: 12
%% degree: 650
%% duplicateId: 0
%% body:
一质量为 $8.00\times10^{4}\Ukg$ 的太空飞船从其飞行轨道返回地面。飞船在离地面高度 $1.60\times10^{5}\Um$ 处以 $7.5\times10^{3}\Ums$ 的速度进入大气层，逐渐减慢至速度为 $100\Ums$ 时下落到地面。取地面为重力势能零点，在飞船下落过程中，重力加速度可视为常量，大小取为 $9.8\Umsq$。（结果保留 2 位有效数字）

\begin{enumerate}
	\item
	分别求出该飞船着地前瞬间的机械能和它进入大气层时的机械能；

	\item
	求飞船从离地面高度 $600\Um$ 处至着地前瞬间的过程中克服阻力所做的功，已知飞船在该处的速度大小是其进入大气层时速度大小的 $2.0\%$。
\end{enumerate}

%% answer:
\jdanswer{
\begin{enumerate}
	\item
	$4.0\times10^{8}\UJ$，$2.4\times10^{12}\UJ$

	\item
	$9.7\times10^{8}\UJ$
\end{enumerate}
}
%% memo:
\memoanswer{
	（1）飞船着地前瞬间的机械能为

	$E_{k0}=\frac{1}{2}mv_{0}^{2}$\ccnum{①}，

	式中，$m$ 和 $v_{0}$ 分别是飞船的质量和着地前瞬间的速率。

	由\ccnum{①}式和题给数据得

	$E_{k0}=4.0\times10^{8}\UJ$\ccnum{②}，

	设地面附近的重力加速度大小为 $g$。飞船进入大气层时的机械能为

	$E_{\text{机}}=\frac{1}{2}mv_{1}^{2}+mgh$\ccnum{③}，

	式中，$v_{1}$ 是飞船在高度 $1.6\times10^{5}\Um$ 处的速度大小。

	由\ccnum{③}式和题给数据得

	$E_{\text{机}}=2.4\times10^{12}\UJ$\ccnum{④}。

	（2）飞船在高度 $h'=600\Um$ 处的机械能为

	$E_{\text{机}}'=\frac{1}{2}m\left(\frac{2.0}{100}v_{1}\right)^{2}+mgh'$\ccnum{⑤}，

	由功能关系得

	$W=E_{\text{机}}'-E_{k0}$\ccnum{⑥}，

	式中，$W$ 是飞船从高度 $600\Um$ 处至着地前瞬间的过程中克服阻力所做的功。由\ccnum{②}\ccnum{⑤}\ccnum{⑥}式和题给数据得

	$W=9.7\times10^{8}\UJ$\ccnum{⑦}。
}

\item
%% number: 25
%% paperName: 2017年全国理综Ⅰ第 25 题 20 分
%% typeId: 6
%% type: 计算题
%% chapter: 电场
%% point: 带电粒子的直线运动
%% method: 无
%% score: 20
%% degree: 350
%% duplicateId: 0
%% body:
真空中存在电场强度大小为 $E_{1}$ 的匀强电场，一带电油滴在该电场中竖直向上做匀速直线运动，速度大小为 $v_{0}$，在油滴处于位置 A 时，将电场强度的大小突然增大到某值，但保持其方向不变。持续一段时间 $t_{1}$ 后，又突然将电场反向，但保持其大小不变；再持续同样一段时间后，油滴运动到 B 点。重力加速度大小为 $g$。

\begin{enumerate}
	\item
	油滴运动到 B 点时的速度；

	\item
	求增大后的电场强度的大小；为保证后来的电场强度比原来的大，试给出相应的 $t_{1}$ 和 $v_{0}$ 应满足的条件。已知不存在电场时，油滴以初速度 $v_{0}$ 做竖直上抛运动的最大高度恰好等于 B、A 两点间距离的两倍。
\end{enumerate}

%% answer:
\jdanswer{
\begin{enumerate}
	\item
	$v_{2}=v_{0}-2gt_{1}$

	\item
	若 B 在 A 点之上，$E_{2}=\left[2-2\frac{v_{0}}{gt_{1}}+\frac{1}{4}\left(\frac{v_{0}}{gt_{1}}\right)^{2}\right]E_{1}$，且 $0<t_{1}<\left(1-\frac{\sqrt{3}}{2}\right)\frac{v_{0}}{g}$ 或 $t_{1}>\left(1+\frac{\sqrt{3}}{2}\right)\frac{v_{0}}{g}$；若 B 在 A 点之下，$E_{2}=\left[2-2\frac{v_{0}}{gt_{1}}-\frac{1}{4}\left(\frac{v_{0}}{gt_{1}}\right)^{2}\right]E_{1}$，且 $t_{1}>\left(\frac{\sqrt{5}}{2}+1\right)\frac{v_{0}}{g}$。
\end{enumerate}
}
%% memo:
\memoanswer{
	（1）设油滴质量和带电荷量分别为 $m$ 和 $q$，油滴速度方向向上为正，油滴在电场强度大小为 $E_{1}$ 的匀强电场中做匀速直线运动，故匀强电场方向向上。在 $t=0$ 时，电场强度突然从 $E_{1}$ 增大至 $E_{2}$ 时，油滴做竖直向上的匀加速运动，加速度方向向上，大小 $a_{1}$ 满足

	$qE_{2}-mg=ma_{1}$\ccnum{①}，

	油滴在 $t_{1}$ 时刻的速度为

	$v_{1}=v_{0}+a_{1}t_{1}$\ccnum{②}，

	电场强度在 $t_{1}$ 时刻突然反向，油滴做匀变速运动，加速度方向向下，大小 $a_{2}$ 满足

	$qE_{2}+mg=ma_{2}$\ccnum{③}，

	油滴在 $t_{2}=2t_{1}$ 时刻的速度为

	$v_{2}=v_{1}-a_{2}t_{1}$\ccnum{④}，

	由\ccnum{①}\ccnum{②}\ccnum{③}\ccnum{④}式得

	$v_{2}=v_{0}-2gt_{1}$\ccnum{⑤}。

	（2）由题意，在 $t=0$ 时刻前有

	$qE_{1}=mg$\ccnum{⑥}，

	油滴从 $t=0$ 到 $t_{1}$ 时刻的位移为

	$s_{1}=v_{0}t_{1}+\frac{1}{2}a_{1}t_{1}^{2}$\ccnum{⑦}，

	油滴在从 $t_{1}$ 时刻到 $t_{2}=2t_{1}$ 时刻的时间间隔内的位移为

	$s_{2}=v_{1}t_{1}-\frac{1}{2}a_{2}t_{1}^{2}$\ccnum{⑧}，

	由题给条件有

	$v_{0}^{2}=2g(2h)$\ccnum{⑨}，

	式中 $h$ 是 B、A 两点之间的距离。

	若 B 点在 A 点之上，依题意有

	$s_{1}+s_{2}=h$\ccnum{⑩}，

	由\ccnum{①}\ccnum{②}\ccnum{③}\ccnum{⑥}\ccnum{⑦}\ccnum{⑧}\ccnum{⑨}\ccnum{⑩}式得

	$E_{2}=\left[2-2\frac{v_{0}}{gt_{1}}+\frac{1}{4}\left(\frac{v_{0}}{gt_{1}}\right)^{2}\right]E_{1}$\ccnum{⑪}，

	为使 $E_{2}>E_{1}$ 应有

	$2-2\frac{v_{0}}{gt_{1}}+\frac{1}{4}\left(\frac{v_{0}}{gt_{1}}\right)^{2}>1$\ccnum{⑫}，

	即当 $0<t_{1}<\left(1-\frac{\sqrt{3}}{2}\right)\frac{v_{0}}{g}$\ccnum{⑬}，

	或 $t_{1}>\left(1+\frac{\sqrt{3}}{2}\right)\frac{v_{0}}{g}$\ccnum{⑭}，

	才是可能的：条件\ccnum{⑬}式和\ccnum{⑭}式分别对应于 $v_{2}>0$ 和 $v_{2}<0$ 两种情形。若 B 点在 A 点之下，依题意有

	$s_{1}+s_{2}=-h$\ccnum{⑮}，

	由\ccnum{①}\ccnum{②}\ccnum{③}\ccnum{⑥}\ccnum{⑦}\ccnum{⑧}\ccnum{⑨}\ccnum{⑮}式得

	$E_{2}=\left[2-2\frac{v_{0}}{gt_{1}}-\frac{1}{4}\left(\frac{v_{0}}{gt_{1}}\right)^{2}\right]E_{1}$\ccnum{⑯}，

	为使 $E_{2}>E_{1}$，应有 $2-2\frac{v_{0}}{gt_{1}}-\frac{1}{4}\left(\frac{v_{0}}{gt_{1}}\right)^{2}>1$\ccnum{⑰}，

	即 $t_{1}>\left(\frac{\sqrt{5}}{2}+1\right)\frac{v_{0}}{g}$\ccnum{⑱}，

	另一解为负，不合题意，已舍去。
}

\item
%% number: 33
%% paperName: 2017年全国理综Ⅰ第 33 题 10 分
%% typeId: 6
%% type: 计算题
%% chapter: 气体性质
%% point: 气体多过程
%% method: 无
%% score: 10
%% degree: 450
%% duplicateId: 0
%% body:
【物理——选修 3-3】（15 分）

（1）（5 分）氧气分子在 $0\UdC$ 和 $100\UdC$ 温度下单位速率间隔的分子数占总分子数的百分比随气体分子速率的变化分别如图中两条曲线所示。下列说法正确的是\xzanswer{ABC}。（填正确答案标号。选对 1 个得 2 分，选对 2 个得 4 分，选对 3 个得 5 分。每选错 1 个扣 3 分，最低得分为 0 分）

\onepicture[h=4.3cm]{\includesvg{figs/33a}}

\fivechoices[answer=ABC]
{图中两条曲线下面积相等}
{图中虚线对应于氧气分子平均动能较小的情形}
{图中实线对应于氧气分子在 $100\UdC$ 时的情形}
{图中曲线给出了任意速率区间的氧气分子数目}
{与 $0\UdC$ 时相比，$100\UdC$ 时氧气分子速率出现在 $0\sim400\Ums$ 区间内的分子数占总分子数的百分比较大}

（2）（10 分）如图，容积均为 $V$ 的汽缸 A、B 下端有细管（容积可忽略）连通，阀门 $K_{2}$ 位于细管的中部，A、B 的顶部各有一阀门 $K_{1}$、$K_{3}$，B 中有一可自由滑动的活塞（质量、体积均可忽略）。初始时，三个阀门均打开，活塞在 B 的底部；关闭 $K_{2}$、$K_{3}$，通过 $K_{1}$ 给汽缸充气，使 A 中气体的压强达到大气压 $p_{0}$ 的 3 倍后关闭 $K_{1}$。已知室温为 $27\UdC$，汽缸导热。

\begin{enumerate}
	\item
	打开 $K_{2}$，求稳定时活塞上方气体的体积和压强；

	\item
	接着打开 $K_{3}$，求稳定时活塞的位置；

	\item
	再缓慢加热汽缸内气体使其温度升高 $20\UdC$，求此时活塞下方气体的压强。
\end{enumerate}

\onepicture[h=4.2cm, align=right]{figs/33b.png}

%% answer:
\jdanswer{
\begin{enumerate}
	\item
	$\frac{V}{2}$，$2p_{0}$

	\item
	活塞位于汽缸 B 顶部

	\item
	$1.6p_{0}$
\end{enumerate}
}
%% memo:
\memoanswer{
	（1）图中曲线下面积为各个速率间隔的分子数占总分子数的百分比之和，即等于 1，故面积相等，故 A 正确；图中虚线对应于氧气分子在 $0\UdC$ 时的情形，其平均动能较小，故 B 正确；图中实线对应于氧气分子在 $100\UdC$ 时的情形，故 C 正确；图中曲线给出了单位速率间隔的分子数占总分子数的百分比，不是分子数目，故 D 错误；与 $0\UdC$ 时相比，$100\UdC$ 时氧气分子速率出现在 $0\sim400\Ums$ 区间内的分子数占总分子数的百分比较小，故 E 错误。

	（2）（i）设打开 $K_{2}$ 后，稳定时活塞上方气体的压强为 $p_{1}$，体积为 $V_{1}$，依题意，被活塞分开的两部分气体都经历等温过程，由玻意耳定律得

	$p_{0}V=p_{1}V_{1}$\ccnum{①}，

	$(3p_{0})V=p_{1}(2V-V_{1})$\ccnum{②}，

	联立\ccnum{①}\ccnum{②}式得 $V_{1}=\frac{V}{2}$\ccnum{③}，

	$p_{1}=2p_{0}$\ccnum{④}。

	（ii）打开 $K_{3}$ 后，由\ccnum{④}式知，活塞必定上升。设在活塞下方气体与 A 中气体的体积之和为 $V_{2}(V_{2}\leq2V)$ 时，活塞下方气体压强为 $p_{2}$，由玻意耳定律得

	$(3p_{0})V=p_{2}V_{2}$\ccnum{⑤}，

	由\ccnum{⑤}式得 $p_{2}=\frac{3V}{V_{2}}p_{0}$\ccnum{⑥}，

	由\ccnum{⑥}式知，打开 $K_{3}$ 后活塞上升直到 B 的顶部为止，此时 $p_{2}'=\frac{3}{2}p_{0}$。

	（iii）设加热后活塞下方气体的压强为 $p_{3}$，气体温度从 $T_{1}=300\UK$ 升高到 $T_{2}=320\UK$ 的等容过程中，由查理定律得

	$\frac{p_{2}}{T_{1}}=\frac{p_{3}}{T_{2}}$\ccnum{⑦}，

	将有关数据代入\ccnum{⑦}式得 $p_{3}=1.6p_{0}$\ccnum{⑧}。
}

\item
%% number: 34
%% paperName: 2017年全国理综Ⅰ第 34 题 10 分
%% typeId: 6
%% type: 计算题
%% chapter: 光学
%% point: 光的折射
%% method: 无
%% score: 10
%% degree: 550
%% duplicateId: 0
%% body:
【物理——选修 3-4】（15 分）

（1）（5 分）如图\ref{2017全国理综Ⅰ34a}，在 $xy$ 平面内有两个沿 $z$ 方向做简谐振动的点波源 $S_{1}(0,4)$ 和 $S_{2}(0,-2)$。两波源的振动图线分别如图\ref{2017全国理综Ⅰ34b}和图\ref{2017全国理综Ⅰ34c}所示，两列波的波速均为 $1.00\Ums$。两列波从波源传播到点 $A(8,-2)$ 的路程差为\tkanswer[1.8cm]{2}$\Um$，两列波引起的点 $B(4,1)$ 处质点的振动相互\tkanswer[1.8cm]{减弱}（填“加强”或“减弱”），点 $C(0,0.5)$ 处质点的振动相互\tkanswer[1.8cm]{加强}（填“加强”或“减弱”）。

\threepicture[labelA=2017全国理综Ⅰ34a, labelB=2017全国理综Ⅰ34b, labelC=2017全国理综Ⅰ34c, numA={（a）}, numB={（b）}, numC={（c）}, hA=2.8cm, hB=2.6cm, hC=2.6cm]
{figs/34a.png}
{figs/34b.png}
{figs/34c.png}

（2）（10 分）如图，一玻璃工件的上半部是半径为 $R$ 的半球体，O 点为球心；下半部是半径为 $R$、高为 $2R$ 的圆柱体，圆柱体底面镀有反射膜。有一平行于中心轴 OC 的光线从半球面射入，该光线与 OC 之间的距离为 $0.6R$。已知最后从半球面射出的光线恰好与入射光线平行（不考虑多次反射）。求该玻璃的折射率。

\onepicture[h=4.4cm, align=right]{figs/34d.png}

%% answer:
\jdanswer{
\begin{enumerate}
	\item
	$2\Um$，减弱，加强

	\item
	$n=\sqrt{2.05}\approx1.43$
\end{enumerate}
}
%% memo:
\memoanswer{
	（1）由勾股定理可知，$S_{1}A=10\Um$，$S_{2}A=8\Um$，路程差为 $2\Um$，由图（b）（c）可知两波源起始振动方向相反，且两波源到 B 点的距离相等，为 $5\Um$，同时传播到 B 点，故 B 处质点的振动相互减弱；由题意可知两列波的波长为 $\lambda=vT=2\Um$，两列波到 C 点的路程差为 $1\Um$，即半个波长，但两列波起始振动方向相反，故 C 处质点的振动相互加强。

	（2）由光路可逆，从半球面射出的光线恰好与入射光线平行，则光路对称，如图所示，由光路的对称性和光路可逆性，与入射光线相对于 OC 轴对称的出射光线一定与入射光线平行。这样，从半球面射入的折射光线，将在圆柱体底面中心 C 点发生反射。

	设光线在半球面的入射角为 $i$，折射角为 $r$，由折射定律得

	$\sin i=n\sin r$\ccnum{①}，

	由正弦定理得

	$\frac{\sin r}{2R}=\frac{\sin(i-r)}{R}$\ccnum{②}，

	由几何关系知，入射点的法线与 OC 间的夹角为 $i$，由题设条件和几何关系得

	$\sin i=\frac{L}{R}$\ccnum{③}，

	式中 $L$ 是入射光线与 OC 间的距离，由\ccnum{②}\ccnum{③}式和题给数据得

	$\sin r=\frac{6}{\sqrt{205}}$\ccnum{④}，

	由\ccnum{①}\ccnum{③}\ccnum{④}式和题给数据得 $n=\sqrt{2.05}\approx1.43$\ccnum{⑤}。

	\onepicture[w=4.6cm]{figs/34d_an.png}
}

\end{enumerate}

\end{document}
